% 第九卷 15.2.7.1 节：分块注意力掩码（正文栏宽，约60 mm）。
% 8个编码步，块大小C=2，左侧保留L=1个历史块。蓝色格为可见的键位置。
\begin{tikzpicture}[x=1mm,y=1mm,figure picture]
  \foreach \i in {0,...,7} {
    \foreach \j in {0,...,7} {
      \pgfmathtruncatemacro{\ci}{floor(\i/2)}
      \pgfmathtruncatemacro{\cj}{floor(\j/2)}
      \pgfmathtruncatemacro{\vis}{and(\cj<=\ci,\cj>=\ci-1)}
      \ifnum\vis=1
        \path[fill=FigureInputFill,draw=FigureInput,line width=.4pt]
          ({5*\j},{-5*\i}) rectangle ++(5,-5);
      \else
        \path[fill=white,draw=FigureGrid,line width=.4pt]
          ({5*\j},{-5*\i}) rectangle ++(5,-5);
      \fi
    }
    \pgfmathtruncatemacro{\k}{\i+1}
    \node[font=\figurefont\scriptsize,anchor=east] at (-1,{-5*\i-2.5}) {$\k$};
    \node[font=\figurefont\scriptsize,anchor=south] at ({5*\i+2.5},1) {$\k$};
  }
  \foreach \k in {0,2,4,6,8} {
    \draw[FigureStroke,line width=.8pt] ({5*\k},0) -- ({5*\k},-40);
    \draw[FigureStroke,line width=.8pt] (0,{-5*\k}) -- (40,{-5*\k});
  }
  \node[font=\figurefont\scriptsize,anchor=south] at (20,6) {键位置$j$};
  \node[font=\figurefont\scriptsize,rotate=90,anchor=south] at (-6,-20) {查询位置$i$};
  \node[font=\figurefont\scriptsize,anchor=west,align=left,text=FigureMuted]
    at (43,-20) {块大小$C=2$\\左侧历史$L=1$块\\粗线为块边界};
\end{tikzpicture}
